% -----------------------------------------------------------------------------
% Q6 -- individual assignment: Hawraa
% Source: individual assignment-hawraa.docx
% Substantive wording is preserved; only document-layout artifacts are removed.
% -----------------------------------------------------------------------------

\definecolor{HawraaNavy}{HTML}{12355B}
\definecolor{HawraaBlue}{HTML}{1769AA}
\definecolor{HawraaTeal}{HTML}{1769AA}
\definecolor{HawraaMint}{HTML}{EDF5FB}
\definecolor{HawraaIce}{HTML}{EEF4F8}
\definecolor{HawraaSand}{HTML}{EDF5FB}

\fancypagestyle{hawraa}{%
  \fancyhf{}%
  \fancyhead[L]{\small\color{AUBGray}BMEN502 \textbar{} Concept Exploration and Testing}%
  \fancyhead[R]{\small\color{HawraaBlue}\bfseries INDIVIDUAL Q6 \textbar{} HAWRAA}%
  \fancyfoot[L]{\small\color{AUBGray}American University of Beirut}%
  \fancyfoot[R]{\small\color{AUBGray}Page \thepage\ of \pageref{LastPage}}%
  \renewcommand{\headrulewidth}{0.7pt}%
  \renewcommand{\headrule}{\hbox to\headwidth{\color{HawraaBlue}\leaders\hrule height \headrulewidth\hfill}}%
  \renewcommand{\footrulewidth}{0pt}%
}

\newcommand{\hawraasectiontitle}[1]{%
  \vspace{0.4em}%
  {\color{HawraaBlue}\Large\bfseries #1}\par
  \vspace{0.25em}%
  {\color{AUBBlack}\rule{\textwidth}{0.8pt}}\par
  \vspace{0.45em}%
}

\newcommand{\hawraaband}[1]{%
  \vspace{0.3em}%
  {\fontsize{13}{15}\selectfont\bfseries\color{HawraaBlue}#1}\par
  \vspace{0.16em}%
  {\color{AUBLine}\rule{\textwidth}{0.55pt}}\par
  \vspace{0.3em}%
}

\newcommand{\hawraacomparison}[1]{%
  \begingroup
  \setlength{\fboxsep}{8pt}%
  \noindent\colorbox{HawraaMint}{%
    \parbox{\dimexpr\linewidth-2\fboxsep-2\fboxrule\relax}{%
      \fontsize{9.15}{11.1}\selectfont #1}}\par
  \endgroup
}

\newcommand{\hawraasource}[1]{%
  \begingroup
  \setlength{\fboxsep}{7pt}%
  \noindent\colorbox{HawraaIce}{%
    \parbox{\dimexpr\linewidth-2\fboxsep\relax}{%
      \fontsize{7.8}{9.25}\selectfont\color{AUBGray}#1}}\par
  \endgroup
}

\newcommand{\hawraacompactband}[1]{%
  \vspace{0.15em}%
  {\fontsize{9.7}{11.4}\selectfont\bfseries\color{HawraaBlue}#1}\par
  \vspace{0.16em}%
  {\color{AUBLine}\rule{\linewidth}{0.55pt}}\par
  \vspace{0.25em}%
}

\newcommand{\hawraacompactcomparison}[1]{%
  \begingroup
  \setlength{\fboxsep}{5pt}%
  \noindent\colorbox{HawraaMint}{%
    \parbox{\dimexpr\linewidth-2\fboxsep-2\fboxrule\relax}{%
      \raggedright\fontsize{6.9}{8.2}\selectfont #1}}\par
  \endgroup
}

\newcommand{\hawraacompactsource}[1]{%
  \begingroup
  \setlength{\fboxsep}{4.5pt}%
  \noindent\colorbox{HawraaIce}{%
    \parbox{\dimexpr\linewidth-2\fboxsep\relax}{%
      \raggedright\fontsize{5.7}{6.7}\selectfont\color{AUBGray}#1}}\par
  \endgroup
}

\newcommand{\hawraachallenge}[1]{%
  \begingroup
  \setlength{\fboxsep}{8pt}%
  \noindent\colorbox{HawraaSand}{%
    \parbox{\dimexpr\linewidth-2\fboxsep-2\fboxrule\relax}{%
      \fontsize{9.1}{11.15}\selectfont #1}}\par
  \endgroup
}

\clearpage
\pagestyle{hawraa}
\thispagestyle{hawraa}
\hypertarget{q6-hawraa}{}

\hawraasectiontitle{Q6. Individual Assignment -- Hawraa}

\begingroup
\setlength{\fboxsep}{9pt}%
\noindent\colorbox{HawraaMint}{%
  \parbox{\dimexpr\linewidth-2\fboxsep\relax}{%
    {\fontsize{9}{11}\selectfont\bfseries\color{HawraaBlue}Question :6}\par
    \vspace{0.2em}
    {\fontsize{12}{14.5}\selectfont\bfseries\color{AUBBlack}
      State-of-the-Art Technologies Comparable to the Proposed Design}\par
  }}\par
\endgroup

\vspace{0.5em}
\noindent
\begin{minipage}[t]{0.486\linewidth}
\vspace{0pt}
\hawraacompactband{1. Össur AeroFit\texttrademark{} Vented Liner--Socket System}

\vspace{0.4em}
{\raggedright\fontsize{7.2}{8.65}\selectfont
The Össur AeroFit\texttrademark{} vented liner--socket system is a commercially
available prosthetic interface designed to reduce moisture accumulation at the
skin--liner interface. The system combines a vented silicone liner containing
an integrated mesh with a vented socket that promotes both vertical and
horizontal airflow during walking. This ventilation mechanism allows
perspiration to escape from the socket, thereby reducing the warm and humid
environment that commonly develops inside conventional prosthetic liners. In a
randomized crossover clinical study involving individuals with transfemoral
amputation, the vented system significantly reduced relative humidity by
approximately 30\% and improved users' perceived sweating without compromising
socket suspension, stability, or overall comfort. (Gnyawali et al. 2023)\par}

\vspace{0.4em}
\hawraacompactcomparison{\textbf{Comparison with the proposed design.} Unlike the
AeroFit\texttrademark{} system, which relies solely on passive ventilation, the
proposed smart prosthetic liner adopts a multifunctional approach. It combines
a reinforced PAM--CMC double-network hydrogel to absorb and manage moisture,
microencapsulated phase-change materials (mPCMs) to regulate temperature, and
embedded flexible sensors to continuously monitor pressure, temperature, and
humidity. By integrating both active microclimate regulation and real-time
physiological monitoring, the proposed design aims not only to reduce moisture
but also to improve thermal comfort and provide objective data for optimizing
socket fit and residual-limb health.}

\vspace{0.38em}
\hawraacompactsource{\textit{Gnyawali, Surya C., Jeffrey A. Denune, Bryce Hockman,
Jóna Valgerður Kristjánsdóttir, Margrét Sól Ragnarsdóttir, Lava R. Timsina,
Subhadip Ghatak, Knut Lechler, Chandan K. Sen, and Sashwati Roy. 2023.
``Moisture Mitigation Using a Vented Liner and a Vented Socket System for
Individuals with Transfemoral Amputation.''} \textit{Scientific Reports}
\textit{ 13: 16557.} \url{https://doi.org/10.1038/s41598-023-43572-2}}
\end{minipage}%
\hfill
\begin{minipage}[t]{0.486\linewidth}
\vspace{0pt}
\begin{center}
\includegraphics[width=\linewidth,height=0.405\textheight,keepaspectratio]
  {assets/q6_hawraa/aerofit_comparison.png}\par
\end{center}
\end{minipage}

% The source DOCX's ``Top of Form'' and ``Bottom of Form'' strings are web-copy
% interface artifacts surrounding the figure and are therefore not typeset.
\vspace{0.55em}
\noindent
\begin{minipage}[t]{0.486\linewidth}
\vspace{0pt}
\hawraacompactband{2. Roliner: Dynamically Adaptive Prosthetic Liner}

\vspace{0.38em}
{\raggedright\fontsize{6.95}{8.35}\selectfont
Roliner is a novel soft metamaterial prosthetic liner designed to dynamically
adapt to changes in the residual limb during daily activities. The liner is
fabricated from silicone elastomers containing embedded millifluidic channels
that can be pneumatically pressurized to modify the liner's shape, volume, and
mechanical properties in real time. Controlled through a portable
electro-pneumatic unit and smartphone application, Roliner enables personalized
adjustment of socket fit, helping accommodate residual-limb volume fluctuations
caused by muscle atrophy, body weight changes, edema, or physical activity.
Preclinical testing demonstrated improved comfort, prosthetic utility,
residual limb health, and fitting while maintaining gait performance comparable
to conventional prosthetic liners.\par}

\vspace{0.38em}
\hawraacompactcomparison{\textbf{Comparison with the proposed design.} While Roliner
focuses on active mechanical adaptation through pneumatic actuation, the
proposed smart prosthetic liner emphasizes passive microclimate regulation
using a reinforced PAM--CMC hydrogel and microencapsulated phase-change
materials (mPCMs), combined with embedded flexible sensors for continuous
monitoring of pressure, temperature, and humidity. This approach avoids the
need for pneumatic hardware while providing thermal regulation, moisture
control, and real-time physiological monitoring within a simpler multilayer
liner architecture.}

\vspace{0.35em}
\hawraacompactsource{Tanriverdi, Ugur, Guglielmo Senesi, Tarek Asfour, Hasan Kurt,
Sabrina L. Smith, Diana Toderita, Joseph Shalhoub, Laura Burgess, Anthony M. J.
Bull, and Firat Güder. 2025. ``Dynamically Adaptive Soft Metamaterial for
Wearable Human--Machine Interfaces.'' \textit{Nature Communications} 16: 2621.
\url{https://doi.org/10.1038/s41467-025-57634-8}}
\end{minipage}%
\hfill
\begin{minipage}[t]{0.486\linewidth}
\vspace{0pt}
\hawraacompactband{3- SmartTemp\textregistered{} Prosthetic Liner (by WillowWood)}

\vspace{0.38em}
{\raggedright\fontsize{6.95}{8.35}\selectfont
The SmartTemp\textregistered{} prosthetic liner is a commercially available
silicone liner designed to improve the thermal environment inside the
prosthetic socket by incorporating phase-change materials (PCMs) into a
conventional silicone liner. PCMs absorb, store, and release thermal energy
during phase transitions, allowing the liner to buffer temperature fluctuations
and delay heat accumulation at the skin--liner interface. In a double-blind
randomized crossover clinical trial involving 16 individuals with transtibial
amputations, the SmartTemp liner significantly reduced residual-limb skin
temperature and perspiration compared with a conventional silicone liner
following 25 minutes of cycling and a 10-minute recovery period. The reduction
in heat and moisture improved the internal socket environment and may help
reduce discomfort, improve suspension, and decrease the risk of skin
complications associated with excessive perspiration.\par}

\vspace{0.38em}
\hawraacompactcomparison{Similar to SmartTemp, the proposed smart prosthetic liner
incorporates microencapsulated phase-change materials (mPCMs) to regulate
residual-limb temperature. However, the proposed design extends this concept by
integrating mPCMs within a reinforced PAM--CMC double-network hydrogel, which
provides simultaneous thermal regulation and passive moisture management. In
addition, embedded flexible sensors continuously monitor pressure, temperature,
and humidity, enabling real-time assessment of the skin--liner interface.
Consequently, the proposed liner offers a multifunctional approach that
combines climate regulation, physiological monitoring, and structural
reinforcement within a single multilayer system rather than relying solely on
PCM-based cooling.}

\vspace{0.35em}
\hawraacompactsource{Wernke, Matthew M., Ryan M. Schroeder, Christopher T. Kelley,
Jeffrey A. Denune, and James M. Colvin. 2015. ``SmartTemp Prosthetic Liner
Significantly Reduces Residual Limb Temperature and Perspiration.''
\textit{Journal of Prosthetics and Orthotics} 27 (4): 134--139.}
\end{minipage}

\clearpage
\thispagestyle{hawraa}
\hawraasectiontitle{Challenges Still Facing the Proposed Design.}

\vspace{0.8em}
\hawraachallenge{Despite integrating several advanced technologies, the
proposed smart prosthetic liner still faces a number of engineering challenges
before clinical implementation. The successful integration of the reinforced
PAM--CMC hydrogel, microencapsulated phase-change materials (mPCMs), embedded
flexible sensors, and the outer structural layer into a single durable and
flexible system remains a major challenge. In addition, the hydrogel must
maintain its mechanical integrity under repeated compression, shear loading,
and prolonged moisture exposure, while the embedded sensors must provide
accurate and stable measurements during continuous bending and daily prosthetic
use. Future development also requires transferring the printed sensors from the
current PET substrate to PDMS to improve flexibility and wearer comfort,
although reliable printing on PDMS is more challenging because of its low
surface energy and reduced ink adhesion. Finally, the complete three-layer
liner has not yet been fabricated or clinically validated, and therefore
comprehensive mechanical, environmental, and user-based testing will be
necessary to evaluate its long-term performance, durability, and effectiveness
in improving prosthetic comfort and residual-limb health.}

\vspace{0.6em}
\hawraachallenge{The development of the proposed smart prosthetic liner
required expertise beyond the knowledge available within our team. One of the
main gaps was advanced biomaterials engineering, particularly in designing and
fabricating reinforced hydrogel systems incorporating microencapsulated
phase-change materials (mPCMs). This expertise was essential to optimize the
hydrogel composition, mechanical durability, moisture absorption, and thermal
regulation required for long-term prosthetic use. A second area of missing
expertise was prosthetics and orthotics, particularly regarding socket design,
liner fabrication, and clinical evaluation of prosthetic interfaces. Such
expertise is necessary to ensure that the proposed liner is compatible with
existing prosthetic systems and meets user comfort and clinical requirements.
Finally, expertise in flexible electronics manufacturing was needed to
integrate printed sensors into a soft wearable substrate such as PDMS while
maintaining reliable electrical performance.}

\vspace{0.6em}
\hawraachallenge{To address these limitations, the design process relied
extensively on published scientific literature, commercial prosthetic
technologies, and expert consultation. Evidence from recent studies was used to
guide the selection of materials and the overall liner architecture, while
commercially available systems such as SmartTemp\textregistered{}, Roliner, and
Össur AeroFit\texttrademark{} were analyzed to identify current design
strategies and limitations. Because of the limited fabrication expertise and
available resources, only the flexible sensor was fabricated as a
proof-of-concept prototype on a PET substrate using a Voltera printer. The
complete three-layer smart liner and sensor integration on PDMS were therefore
proposed as future work requiring collaboration with specialists in
biomaterials, prosthetics, and flexible electronics.}

\clearpage
\thispagestyle{hawraa}
\vspace*{\fill}
\begin{center}
\includegraphics[width=\linewidth,height=0.78\textheight,keepaspectratio]
  {assets/q6_hawraa/three_layer_liner_design.png}\par
\end{center}
\vspace*{\fill}

\clearpage
\pagestyle{fancy}
\fancyhf{}
\fancyhead[L]{\small\color{AUBGray}BMEN502 \textbar{} Concept Exploration and Testing}
\fancyhead[R]{\small\color{AUBRed}\bfseries SMART LINER}
\fancyfoot[L]{\small\color{AUBGray}American University of Beirut}
\fancyfoot[R]{\small\color{AUBGray}Page \thepage\ of \pageref{LastPage}}
\renewcommand{\headrulewidth}{0.6pt}
\renewcommand{\headrule}{\hbox to\headwidth{\color{AUBRed}\leaders\hrule height \headrulewidth\hfill}}
\renewcommand{\footrulewidth}{0pt}
